\section{韩：小国的官僚化实验}

韩国的地理位置决定了它很难像齐、楚那样依靠广阔腹地，也很难像秦那样依靠关中屏障。它夹在魏、赵、秦、楚等强国之间，任何一次战争都可能触到国都附近。对这样的国家而言，资源不多，便更需要知道官吏是否尽职、军队能否迅速调动、命令能否穿过贵族关系抵达地方。

\begin{event}{公元前375年}{韩灭郑}{可靠纪年}{郑地}\eventanchor{WQ-0375-HAN}{战国·韩灭郑}
\term{这一年发生了什么}：韩国灭郑，获得向东扩展的土地与人口。

\term{后来改变了什么}：这使韩从晋国旧卿族的一支，逐渐成为拥有较完整领土的战国国家；但新得的土地也意味着更高的行政、军事和财政需求。小国的扩张并不会消除小国的不安全，反而常把它推到更强邻国的视线中。
\end{event}

\begin{event}{约公元前354年至前337年}{申不害相韩}{王年较可靠，细节有争议}{韩}\eventanchor{WQ-0354-SHENBUHAI}{战国·申不害相韩}
\term{这一阶段发生了什么}：申不害长期辅佐韩昭侯。后世将他的政治实践概括为“术”：通过形名、考核和君主对官吏的驾驭，使有限的国家资源不被臣下各自占有。

\term{事情为什么走到这里}：韩国不能只靠封邑贵族之间的默契维持国家。面对强邻，它需要更可计算的任官和监督机制；所谓“术”，首先是小国在高压竞争中试图让官僚机器少出错的一种技术。

\term{后来改变了什么}：韩国一度“国治兵强”的后世评价，至少说明行政整顿被看作国家竞争的重要部分。但行政技术不能改变地缘位置本身；韩最终仍是秦东进最先面对的国家之一。
\end{event}

\begin{sourcenote}[申不害的“术”是谁保存下来的]
《申子》原本已佚，现存说法多来自后世类书和《韩非子》的概括与批评。韩非的解释极重要，却有自己的理论目的。这里把“形名”“考核”和君主驾驭官吏视作理解韩国行政技术的线索，不把后代哲学分类倒灌成申不害本人完整不变的体系。
\end{sourcenote}

\begin{turningpoint}[制度得失：精细行政能抵消多少地缘压力]
申不害式的行政技术可以减少信息失真、约束官吏和集中有限资源，这是韩国的所得；它不能凭空创造纵深、人口和战略缓冲，这是韩国的限制。战国国家竞争让制度变得越来越重要，也让制度的边界变得越来越清楚。
\end{turningpoint}
